\documentclass[conference]{IEEEtran}
\usepackage[utf8]{inputenc}
\usepackage{amsmath,amsfonts}
\usepackage{graphicx}
\usepackage{listings}
\usepackage{booktabs}
\usepackage{multirow}
\usepackage[linesnumbered,ruled,vlined]{algorithm2e}

\lstdefinelanguage{PDDL}{
	morekeywords={define,domain,problem,:domain,:init,:goal,:action,:parameters,:precondition,:effect,:objects,:predicates,:functions,:types,:requirements,:metric},
	sensitive=false,
	morecomment=[l]{;},
	morestring=[b]",
}
\lstset{
	language=PDDL,
	basicstyle=\ttfamily\small,
	breaklines=true,
	columns=flexible,
	frame=single,
	captionpos=b
}

\large

\title{A BDI Agent for Pickup-and-Delivery Tasks in Dynamic Environments}
\author{
	\IEEEauthorblockN{Mauro Antonio de Palma - 256175}
	\IEEEauthorblockA{Email: mauroantonio.depalma@studenti.unitn.it}
	\and
	\IEEEauthorblockN{Marco Crespi - Matricola}
	\IEEEauthorblockA{Email: marco.crespi@studenti.unitn.it}
	
}
	
\begin{document}
\maketitle


\section{Introduction}

Autonomous agents operating in dynamic environments must balance incomplete knowledge, evolving goals, and timely decision-making. In delivery scenarios—such as logistics, warehouse automation, or simulated pickup-and-delivery tasks—agents are often required to identify high-value items, navigate to their locations, and transport them to designated destinations under constraints such as time pressure or partial observability.

This paper presents a Belief-Desire-Intention (BDI) agent designed for a pickup-and-delivery task in a grid-based environment. The agent begins with limited information: it knows the delivery location, which tiles are navigable, and which can potentially spawn parcels. However, the locations of the parcels themselves are initially unknown and must be discovered through exploration. Furthermore, each parcel's score decays over time, which incentivizes the agent to prioritize its goals efficiently.

By modeling decision-making through the BDI architecture, the agent can dynamically update its beliefs based on observations, generate context-dependent desires, and commit to intentions that balance exploration and exploitation. This work focuses on how the agent selects and pursues intentions under uncertainty, with the aim of maximizing overall delivery performance.

Beyond the single-agent scenario, the DeliverooJS environment also supports collaborative settings via an inter-agent messaging system. We extend our agent design to leverage this capability, allowing agents to coordinate, share perceptions, and optimize parcel collection as a team. This dual focus allows us to explore both individual reasoning under uncertainty and collective planning via lightweight communication.

To evaluate our approach, we introduce a benchmarking suite that compares our BDI agent with other competitive strategies in both single-agent and collaborative scenarios. These benchmarks help assess not only raw performance (e.g., delivery score, responsiveness), but also adaptability and coordination effectiveness in multi-agent settings. By varying map layouts, spawn frequencies, and team sizes, we demonstrate the robustness and flexibility of our agent design across diverse conditions.


\section{Agent Architecture}

The agent in this work is designed using the Belief-Desire-Intention (BDI) model, a well-established framework for building rational agents that operate in dynamic, partially observable environments. In this section, we detail how the BDI paradigm is instantiated for our pickup-and-delivery agent within the DeliverooJS game.


\subsection{Beliefs}

Beliefs represent the agent’s internal state of knowledge about the world. In our implementation, the belief base includes:
\begin{itemize}
	\item The static map layout (e.g., navigable tiles, parcel spawn points, delivery location),
	\item The agent’s own position,
	\item Opponents agent's positions
	\item Observed parcel locations and their current scores (when visible),
	\item Messages received from other agents (e.g., shared discoveries),
	\item Known teammate positions (if communication is enabled).
\end{itemize}

Beliefs are updated through direct perception (sensor input from the DeliverooJS API) and inter-agent communication. The belief update function fuses new perceptions into the existing knowledge base, overwriting outdated observations or appending new information as needed.

\subsection{Desires}

Desires capture the agent’s possible goals, given its current beliefs. In our context, the agent dynamically constructs a set of candidate desires such as:
\begin{itemize}
	\item Deliver the parcel at position $p$ before its value decays significantly,
	\item Explore an unexplored tile to find potential parcels,
	\item Move toward a teammate to coordinate a pickup,
	\item Avoid congestion in high-traffic zones.
\end{itemize}

Each desire is associated with a utility score based on criteria like proximity, parcel value decay, and current task urgency. For instance, a parcel with high value and short delivery distance is prioritized over distant or low-value ones.

\subsection{Intentions and Planning}

Intentions are the subset of desires that the agent commits to pursuing. Commitment is governed by a filtering function that considers:
\begin{itemize}
	\item The utility of each desire,
	\item Team-level priorities (if operating in collaborative mode).
\end{itemize}

Once intentions are selected, the agent invokes a planning function to generate an executable intention plan $\pi$ — a sequence of actions such as \texttt{move}, \texttt{pickup}, and \texttt{deliver}. The plan is constructed using the goal-directed search A*, constrained by the current beliefs.

\subsection{Communication and Coordination}

In collaborative scenarios, agents utilize DeliverooJS’s messaging system to share critical beliefs. This includes discovered parcel locations, unreachable areas, or requests for help (e.g., \textit{``You're in a better position to deliver package p1. Let's meet at (5,3) so I can pass it to you''}). The belief base is extended to include received messages, which influence both desire generation and intention filtering.

Communication is kept lightweight to minimize overhead and ensure reactivity. Messages are structured and timestamped, and each agent processes them before the next BDI cycle.

\subsection{BDI Control Loop}

The complete reasoning cycle of the agent is described in Algorithm~\ref{alg:bdi-loop}. This loop ensures the agent continuously adapts its plans based on environmental changes and new information.

\begin{algorithm}
	\caption{BDI Control Loop}
	\label{alg:bdi-loop}
	\KwIn{Initial beliefs $B_0$, desire evaluation function $\texttt{GenerateDesires}$, plan library $\mathcal{P}$}
	\KwOut{Action stream sent to actuators}
	$B \gets B_0$\;
	\While{agent is active}{
		$O \gets \texttt{PerceiveEnvironment()}$\;
		$M \gets \texttt{ReceiveMessages()}$\;
		
		$B \gets \texttt{UpdateBeliefs}(B, O, M)$\;
		
		$D \gets \texttt{GenerateDesires}(B)$\;
		$I \gets \texttt{FilterIntentions}(D, B)$\;
		
		\If{$\texttt{NeedNewPlan}(I, B)$}{
			$\pi \gets \texttt{Plan}(I, B, \mathcal{P})$\;
			$\texttt{CommitTo}(\pi)$\;
		}
		
		$\alpha \gets \texttt{ExecuteNextAction}(\pi)$\;
		\texttt{SendToActuators}($\alpha$)\;
	}
\end{algorithm}

\section{Communication protocol}

The DeliverooJS environment enables agents to exchange information using two primitive messaging methods:
\begin{itemize}
	\item \texttt{shout(message)} — broadcasts a message to all agents within communication range,
	\item \texttt{say(agentId, message)} — sends a direct message to a specific agent identified by its unique ID.
\end{itemize}

At runtime, each agent periodically broadcasts a \texttt{HELLO} message using \texttt{shout}. This behavior is crucial in dynamic scenarios where agents can spawn at any time during execution. Broadcasting \texttt{HELLO} ensures that newly spawned agents can discover existing teammates and integrate into a collaborative strategy.

Once an agent receives a valid \texttt{HELLO} message and confirms it matches expected grammar and cryptographic properties, it adds the sender to its list of known allies. At that point, the agent can initiate collaborative behaviors, such as sharing parcel discoveries, coordinating pickup assignments, or load balancing deliveries.

\subsubsection{Security and Trust}

To ensure message confidentiality and prevent man-in-the-middle attacks, all inter-agent messages are encrypted using the AES (Advanced Encryption Standard) symmetric-key algorithm. Each agent shares a pre-established key known only to legitimate team members.

In addition to encryption, message integrity and authenticity are checked using an internal message grammar. Agents are designed to only accept messages that conform to this strict format. If a received message fails grammatical validation, it is discarded as potentially spoofed or corrupted.

This two-layer approach—cryptographic encryption combined with grammatical validation—helps establish trust and resilience in adversarial or noisy communication settings, where both misinformation and interference are possible.

[\textit{Detailed message formats, encryption keys, and validation schemes will be formalized in a dedicated appendix or protocol specification section.}]

\section{Beliefs}

The belief base defines the agent’s internal model of the environment and is updated over time based on direct perceptions and messages from allies. It includes both static world knowledge and dynamic runtime observations.

\subsection{Static Map Knowledge}

At initialization, the agent is provided with full knowledge of the environment's layout. Each tile in the grid map is labeled using the following enumeration:

\begin{itemize}
	\item \texttt{WALKABLE}: tiles the agent can navigate through,
	\item \texttt{SPAWN}: tiles where parcels may appear,
	\item \texttt{DELIVERY}: the known drop-off location for parcels,
	\item \texttt{NON\_WALKABLE}: obstacles or unreachable areas.
\end{itemize}

Before its main decision loop begins, the agent preprocesses the tile map into a graph structure, where nodes represent \texttt{WALKABLE} tiles, and edges link adjacent nodes (using 4-directional movement). This graph supports efficient path planning using the A* algorithm. 

In addition, the agent analyzes the graph’s connected components to identify unreachable sections. This ensures it does not attempt to reach parcels or tiles in disconnected areas of the map.

\subsection{Pathfinding and Heuristics}

Navigation relies on A* search, where the cost function combines path length and a heuristic estimate of remaining distance. The heuristic used is the Manhattan distance:

\[
h(n) = |x_n - x_g| + |y_n - y_g|
\]

where $(x_n, y_n)$ is the current tile and $(x_g, y_g)$ is the goal. This reflects the environment’s movement constraints and ensures consistent, admissible guidance for the planner.

\subsection{Dynamic Beliefs}

Over time, the agent updates its beliefs to reflect changing or discovered information. Dynamic beliefs include:

\begin{itemize}
	\item \textbf{Parcel positions and scores:} As parcels appear within the agent’s field of view, or are communicated by allies, their location and time-sensitive score are stored. Since scores decay over time, the agent continuously updates or prunes this data.
	
	\item \textbf{Opponent positions:} When visible, opponents’ positions are recorded. These beliefs help avoid conflicts (e.g., when racing to the same parcel) or to identify blocked routes.
	
	\item \textbf{Ally intentions (collaborative mode):} In team scenarios, allies may share information about discovered parcels and opponents. These shared beliefs are integrated to coordinate pickup strategies and reduce redundancy. Special mention goes to the \textbf{Hand-off} protocol, which facilitates collaboration in cases of blocked paths, as well as the map \textbf{exploration split} strategy to optimize area coverage.

	
	\item \textbf{Exploration state:} The agent keeps track of visited and unvisited tiles, allowing it to prioritize unexplored areas when no known parcels are available.
\end{itemize}

All dynamic beliefs are time-sensitive. Data is associated with timestamps or observation turns and is discarded after expiration or when proven obsolete (e.g., a parcel has already been picked up).

\subsection{Exploration Heuristics}

When no known parcels are immediately reachable, the agent engages in active exploration. Its goal is to discover new parcels by visiting promising areas (\texttt{SPAWN} tiles) where parcels are likely to appear.

To prioritize unexplored regions, the agent assigns a dynamic score to each spawn tile using a multi-factor exploration heuristic. This heuristic favors tiles that are:

\begin{itemize}
	\item Less frequently visited,
	\item Closer to the agent's current position,
	\item Less crowded by other agents,
	\item Slightly randomized to avoid deterministic loops.
\end{itemize}

Formally, the exploration score $S(p)$ for a tile $p$ at position $\mathbf{p}$ is defined as:

\[
S(p) = \frac{1}{v(\mathbf{p}) + 1} - 0.05 \cdot d(\mathbf{p}) - a(\mathbf{p}) + r
\]

where:
\begin{itemize}
	\item $v(\mathbf{p})$ is the number of times the agent has visited tile $\mathbf{p}$,
	\item $d(\mathbf{p})$ is the Manhattan distance from the agent’s current position to $\mathbf{p}$,
	\item $a(\mathbf{p})$ is the estimated density of other agents near $\mathbf{p}$,
	\item $r$ is a small random factor sampled uniformly in $[0, 1)$ to break ties and encourage stochasticity.
\end{itemize}

The agent maintains a priority queue of candidate tiles ranked by $S(p)$ and selects the top-scoring tile for exploration. This mechanism allows for efficient coverage of the map and adapts to both teammate activity and environmental congestion.

%! Desires subsection
\section{Desires}

Desires represent the goals that the agent might pursue, based on its current beliefs. In our system, each desire is a structured object defined by a type, priority value, optional target position, and additional contextual information. Desires are generated dynamically in each BDI reasoning cycle and are later filtered to determine which intention(s) to commit to.

\subsection{Desire Representation}

Each desire belongs to one of several predefined types, formalized in the \texttt{DesireTypes} enumeration:

\begin{itemize}
    \item \texttt{DELIVER\_PARCEL}: Bring one or more carried parcels to the delivery zone,
    \item \texttt{PICKUP\_PARCEL}: Navigate to a known parcel and pick it up,
    \item \texttt{EXPLORE\_ENVIRONMENT}: Investigate unexplored spawn tiles,
    \item \texttt{PUT\_DOWN\_PARCEL}: Drop parcels on the delivery point (e.g., after a handoff),
    \item \texttt{PICKUP\_HANDOFF}: Navigate to a known parcel and ask for help delivering it
\end{itemize}

Each desire is also associated with a numeric priority from the \texttt{DesirePriorities} enumeration. These priorities determine the desirability of a goal and influence which intentions are selected. Examples include:

\begin{itemize}
    \item \texttt{PRIORITY\_PICKUP} = 100 (e.g., urgent parcel pickup),
    \item \texttt{PRIORITY\_DELIVERY} = 100,
    \item \texttt{DELIVERY} = 90,
    \item \texttt{HANDOFF\_PRIORITY} = 85,
    \item \texttt{PICKUP} = 80 (standard pickups),
    \item \texttt{EXPLORATION} = 50 (spawn tile exploration).
\end{itemize}

Each BDI cycle triggers a regeneration of the desire set. As new parcels are discovered, agent roles evolve, or tiles are explored, old desires are discarded and new ones are created. This dynamic refresh allows the agent to stay responsive and adaptive to environmental changes.

\subsection{Optional Planning Mode (PDDL)}

In addition to the default reactive behavior, this agent supports an optional planning mode based on PDDL (Planning Domain Definition Language) \cite{mcdermott1998pddl}. This mode is intended for users who want the agent to reason over longer-term sequences of actions using classical planning techniques.

\paragraph{Motivation.}
The PDDL mode is particularly useful in structured delivery scenarios where multiple packages must be picked up and delivered with consideration for path costs and decaying rewards. Rather than making short-sighted decisions at every step, the planner computes a full plan that aims to optimize overall performance based on the current world state.

\paragraph{Planning Structure.}
The domain defines three core actions: \texttt{move}, \texttt{pick-up}, and \texttt{deliver}. Each action contributes to changes in the agent's state, the spatial distribution of parcels, and the evolution of scoring or cost-related fluents. The planning objective is to minimize a total cost function that captures both movement effort and penalties associated with parcel score decay.

The cost of executing a \texttt{move} action from location \( A \) to location \( B \) is computed as:

\[
\text{cost}_{\text{move}} = \text{ManhattanDistance}(A, B) \times \left( \frac{\text{moveDuration}}{\text{scoreDecayInterval}} \right)
\]

This formulation reflects the trade-off between spatial traversal and temporal decay, where longer movements incur proportionally higher penalties due to the effect of delayed delivery on parcel score.



The corresponding problem file encodes the concrete instance: one agent located at a specific position, several packages placed across the map, and a known delivery location. Each parcel has an associated negative score (representing urgency), and distances between all relevant locations are defined using metric functions.

\paragraph{Core Logic.}
The planner's goal is to ensure that the agent has delivered at least one parcel and is no longer carrying any. This is formalized in the goal clause as:
\begin{verbatim}
(:goal (and
  (>= (delivered-parcels) 1)
  (= (carrying-parcels) 0)
))
\end{verbatim}

To reach this state, the planner chooses a sequence of move, pickup, and deliver actions that jointly satisfy the goal while minimizing the total cost, which integrates both physical distance and the negative reward associated with late deliveries.

\paragraph{Composite Cost Model.}
Unlike typical planning domains focused only on distance, the agent's domain includes custom cost modeling:
\begin{itemize}
    \item \texttt{move} actions increase \texttt{total-cost} by the distance between locations.
    \item \texttt{deliver} actions decrease \texttt{total-cost} by the (negative) score of the delivered package.
\end{itemize}

As a result, the agent is incentivized to prioritize high-value parcels even if they are farther away, as their negative score contributes to reducing the final cost function.

\paragraph{Optional Use and Flexibility.}
The PDDL planning mode is optional and disabled by default. Users can enable it through a configuration flag. When activated, the agent queries an external planner with the current domain and problem instance and receives a sequence of grounded actions to execute. Given the temporal and numeric characteristics of the domain, we chose \textbf{ENHSP} (Expressive Numeric Heuristic Search Planner) as the backend planner. \textbf{ENHSP} supports PDDL domains with numeric fluents, temporal constraints, and action costs, which are essential to modeling parcel decay, movement penalties, and delivery rewards in our environment. Its support for heuristic-guided search with cost-sensitive objectives makes it particularly well-suited to compute high-quality plans under constrained conditions.


This approach is ideal in deterministic settings or when a user needs a globally optimized delivery plan. However, due to its offline nature and computational overhead, it may be less suited to rapidly changing environments where new parcels or obstacles appear dynamically.

To illustrate how the planning domain models these dynamics, we provide below a representative excerpt of the PDDL \texttt{domain} and \texttt{problem} definitions used by the agent when invoking the external planner.


\begin{lstlisting}[language=PDDL, caption={Example of generated PDDL domain and problem showing delivery planning logic.}]
(define (domain delivery)
  (:requirements :strips :typing :action-costs)
  (:types agent package location)

  (:predicates
    (at ?a - agent ?l - location)
    (at-pkg ?p - package ?l - location)
    (carrying ?a - agent ?p - package)
    (is-delivery ?l - location)
    (available ?p - package)
    (different ?from ?to - location)
  )

  (:functions
	(score ?p - package)
	(distance ?from ?to - location)

	(total-cost)
	(carrying-parcels)
	(delivered-parcels)
  )
  
  (:action move
  	:parameters (?a - agent ?from - location ?to - location)
  	:precondition (and
  	  (at ?a ?from)
  	  (different ?from ?to)
  	)
  	:effect (and
  	  (not (at ?a ?from))
  	  (at ?a ?to)
  	  (increase (total-cost) (distance ?from ?to))
  	)
  )
  
  (:action pick-up
  	:parameters (?a - agent ?p - package ?l - location)
  	:precondition (and
	  (at ?a ?l)
  	  (at-pkg ?p ?l)
  	  (available ?p)
  	)
  	:effect (and
  	  (carrying ?a ?p)
  	  (not (at-pkg ?p ?l))
 	  (increase (carrying-parcels) 1)
  	)
  )
  

  (:action deliver
    :parameters (?a - agent ?p - package ?l - location)
    :precondition (and
      (at ?a ?l)
      (carrying ?a ?p)
      (is-delivery ?l)
    )
    :effect (and
      (not (carrying ?a ?p))
      (decrease (carrying-parcels) 1)
	  (decrease (total-cost) (score ?p))
	  (increase (delivered-parcels) 1)
    )
  )
)

(define (problem delivery-instance)
  (:domain delivery)
  (:objects
    agent1 - agent
    pkg1 - package
    locA locB locC - location
  )

  (:init
    (at agent1 locA)
    (at-pkg pkg1 locB)
    (is-delivery locC)
    (available pkg1)
    (different locA locB)
    (different locB locC)
    (= (score pkg1) -20)
    (= (distance locA locB) 2)
    (= (distance locB locC) 3)
    (= (total-cost) 0)
	(= (carrying-parcels) 0)
	(= (delivered-parcels) 0)
  )

(:goal (and
	(>= (delivered-parcels) 1)
	(= (carrying-parcels) 0)
))

  (:metric minimize (+ (total-cost) (* -20 (delivered-parcels))))

)
\end{lstlisting}


%! Intentions subsection
\section{Intentions}

\subsection{Definition and Role}

In the Deliveroo agent, intentions represent the agent’s chosen commitments: the plans it is actively pursuing. While desires are possible goals the agent might want to achieve, intentions are selected desires that have been turned into concrete actions. These intentions guide the agent’s behavior over time and are at the core of its deliberative process.

\subsection{Execution Loop}

At every decision step, the agent follows a simple control loop: it selects the most promising desire, translates it into an intention, and attempts to execute that intention. If the intention is successful and reaches its goal, it is removed. If not, it may be retried or eventually discarded. This process is illustrated below:

\begin{algorithm}
	\DontPrintSemicolon
	\caption{Deliberation and Intention Execution Loop}
	\While{agent is active}{
		$desire \gets$ select highest-priority desire\;
		$intention \gets$ generate intention from $desire$\;
		execute($intention$)\;
		\If{intention is complete}{
			remove $intention$ from memory\;
		}
		\ElseIf{intention repeatedly fails}{
			mark $desire$ as failed and discard $intention$\;
		}
	}
\end{algorithm}

This simple but flexible loop enables the agent to dynamically adapt its behavior as goals and circumstances evolve in the environment.

\subsection{Composite Intentions for Pickup and Delivery}

A noteworthy case occurs with parcel pickup and delivery tasks. These usually require a **sequence** of actions: first moving to the parcel location, and then picking it up. For this reason, the agent generates two separate intentions:
\begin{enumerate}
	\item A \texttt{MOVE} intention to navigate to the parcel location.
	\item A \texttt{PICK\_UP} intention to collect the parcel once the agent arrives.
\end{enumerate}

This modular decomposition allows the agent to treat complex goals as a series of manageable steps, improving robustness and reactivity. For instance, if the MOVE intention fails due to a blocked path, it can be abandoned or rerouted without affecting the overall pickup goal.

The same logic applies to deliveries: the agent first moves toward the delivery point, and only once it arrives does it execute the actual drop-off action. This separation of concerns allows the system to better handle interruptions, rerouting, or replanning based on updated environmental information.



%! Handoff subsection
\section{Handoff}

Handoff is a process where one agent requests assistance from another agent to facilitate the delivery of parcels.

This situation occurs when the agent carrying the parcels recognizes that another agent is in a better position to deliver the package more efficiently.
This may be due to being closer to a drop box or because the agent currently carrying the parcels is unable to reach a drop box.
By coordinating handoff, agents can optimize the delivery process, ensuring that parcels are delivered efficiently.

\subsection{Coordination}

Coordination between agents is performed by exchanging messages between them.

When an agent determines that a handoff is beneficial, it sends a request to another agent.
An agent evaluates a handoff as beneficial when the other agent is closer to a delivery position.

When an agent sends a handoff request, it calculate a meeting point to include in the request.
This meeting point is determined by finding the average position between the agent making the handoff request and the agent receiving it.
The agent sending the request immediately begins moving toward the meeting point.

The agent receiving the handoff request evaluates whether it is able to reach the location indicated as the meeting point and responds to the request positively or negatively based on this evaluation.

\subsection{Performing the handoff}

When a handoff request is accepted, the agent that created the request deposits the parcels it carries at the designated meeting point.
Once deposited, the parcels will not be considered in future pickup evaluations of the agent.

Instead, the responding agent moves to the handoff location where the requesting agent deposits the parcels and proceeds to collect them in a standard pickup interaction.

At this point, the handoff is considered complete by both agents.


%! Experimental results subsection
\section{Experimental Results}

To evaluate the performance of the Deliveroo agent, we conducted experiments on a benchmark suite of 10 delivery environments: 5 designed for single-agent scenarios and 5 for collaborative (multi-agent) scenarios. Each configuration was run for 5 minutes under the same computational constraints. The agent's score is defined by the total value of delivered parcels, accounting for decay over time and movement cost.

\subsection{Evaluation Protocol}

For each level, the agent was initialized in a fixed position and allowed to interact with its environment using either its default deliberative mode or the optional PDDL planner. The evaluation metric is the total score achieved at the end of the 5-minute episode.

\subsection{Score Results}

Table~\ref{tab:results} summarizes the scores obtained in each level.

\begin{table}
	\centering
	\caption{Performance per level: total score, average score per parcel, and number of delivered parcels for PDDL and non-PDDL agents}
	\label{tab:results}
	\begin{tabular}{|l|ccc|ccc|}
		\toprule
		\textbf{Level} & \multicolumn{3}{c|}{\textbf{With PDDL}} & \multicolumn{3}{c|}{\textbf{Without PDDL}} \\
		\cmidrule{2-7}
		& \textbf{Total} & \textbf{Avg} & \textbf{\# Parcels} & \textbf{Total} & \textbf{Avg} & \textbf{\# Parcels} \\
		\midrule
		25c1\_1 & 330 & 10.0 & 33 & 380 & 10.0 & 38 \\
		25c1\_2 & 222 & 6.59 & 32 & 1079 & 19.74 & 53 \\
		25c1\_3 & 1104 & 13.51 & 79 & 1303 & 14.46 & 87 \\
		25c1\_4 & 1169 & 67.18 & 18 & 2569 & 77.85 & 33 \\
		25c1\_5 & 2616 & 61.57 & 39 & 3955 & 68.82 & 57 \\
		25c1\_6 & 619 & 20.64 & 28 & 839 & 22.75 & 36 \\
		\midrule
		% optional average row
	\end{tabular}
\end{table}



\subsection{Discussion}

As shown in Table~\ref{tab:results}, the collaborative mode consistently outperforms the single-agent setup. This improvement stems from the agents' ability to share parcel and opponent beliefs, split the map for exploration, and avoid redundant pickups. Notably, the gain is most significant in larger or more spatially distributed levels, where coordination reduces travel overhead and improves delivery rate.

These results validate the effectiveness of the BDI architecture in both standalone and cooperative contexts, and they demonstrate that structured agent collaboration can yield significant efficiency gains in decentralized delivery tasks.


\bibliographystyle{IEEEtran}
\bibliography{references}

\end{document}